Este capítulo recoge los resultados obtenidos con el sistema desarrollado, abarcando tanto la validación del simulador frente a los datos
operativos reales como la experimentación mediante modificación paramétrica de los escenarios. Para ello, se estructura en cuatro secciones.
La primera, a modo introductorio, presenta el funcionamiento operativo del sistema completo. La segunda presenta los distintos indicadores de
rendimiento (KPIs) considerados y su visualización para cada solución. La tercera aborda la validación operativa del simulador y la comparativa
entre las dos soluciones presentadas en el capítulo anterior. Por último, la cuarta cubre la experimentación mediante modificación paramétrica
del sistema, contrastando los resultados obtenidos frente a las hipótesis planteadas.

\section{Funcionamiento operativo del sistema final}\label{sec:general_behaviour}
Si bien el núcleo del presente trabajo lo constituye el simulador multiagente desarrollado sobre GAMA, su explotación práctica requiere de un medio
que permita a un usuario no especializado configurar, lanzar y consultar simulaciones sin interactuar directamente con la plataforma.
Con este propósito, y conforme a los requisitos definidos en el marco de la colaboración con la EMT~(recogidos en el anexo~\ref{anex:requirements}),
el sistema final incorpora una interfaz web que actúa como capa intermedia entre el usuario y el simulador. Cabe destacar que dicha interfaz no 
constituye un objetivo de investigación en sí misma, sino el instrumento mediante el cual se orquestan los lanzamientos y se visualizan los 
resultados que se presentan a lo largo de este capítulo.

La interfaz centraliza tres funciones: la parametrización de los escenarios de simulación, el lanzamiento y almacenamiento de las ejecuciones,
y la visualización de los indicadores de rendimiento resultantes. Con el fin de garantizar la reproducibilidad del despliegue y eliminar
dependencias del sistema anfitrión, el conjunto completo se distribuye sobre una imagen de \textit{Docker}, de modo que simulador, interfaz y
persistencia de resultados se ejecutan de forma autocontenida.

Los resultados presentados en las secciones posteriores derivan tanto de figuras como de instantáneas de dicha interfaz. La exposición acerca
de cómo se visualiza cada indicador de forma individual se presenta a continuación en la siguiente sección.

\section{Estructura y visualización de los KPIs}\label{sec:kpi_visualization}

Manteniendo la estructura de métricas presentada en la sección~\ref{sec:kpi_extraction}, se presenta para cada categoría su
correspondiente indicador mostrado en la interfaz. Cabe destacar que tanto las figuras generadas como los valores relativos a cada 
indicador se computan una única vez al finalizar cada simulación, y no bajo demanda en el momento de la consulta. Este cómputo se 
realiza para todas las combinaciones de filtrado posibles, por lo que el usuario final experimentará tiempos de respuesta prácticamente
inmediatos en el muestreo de resultados. 

De forma previa a la presentación de cada métrica, resulta imperioso remarcar que, con el objetivo de sintetizar alertas visuales
para la detección de posibles puntos críticos, donde se estén dando resultados anómalos, se ha dispuesto de un sistema heurístico basado
en umbrales que, mediante clasificación por colores, permita al usuario identificar dichos puntos para una rápida exploración de estos.
Dicho sistema resulta implementado tanto en las \textbf{métricas a nivel de ruta} como en las \textbf{métricas a nivel de autobús}. En cada apartado
correspondiente se expande la información relevante sobre los umbrales determinados.









\subsection{Métricas a nivel general}

\subsection{Métricas a nivel de ruta}

\subsection{Métricas a nivel de autobús}


% \subsection{Sistema de catalogación de resultados}
% Resulta importante exponer el \textbf{criterio de identificación de resultados destacables} implementado para la interfaz.

% Este consiste en una clasificación por colores según distintos parámetros, que, a modo de nivel de alerta visual, permiten identificar
% posibles puntos donde se estén dando resultados anómalos. Para ello, se especifican valores a modo de frontera para cada uno de los cuatro
% campos a estudiar.

% Concretamente:

% \begin{enumerate}
%     \item \textit{Valor de regularidad}:

%     Se establece el valor~\textbf{80} como valor mínimo requerido para no activar este indicador.

%     \item \textit{Porcentaje de simulaciones con retraso}:
    
%     Se activa este indicador siempre que el indicador relativo a este campo para cada expedición de un un autobús muestre un porcentaje
%     de réplicas con retraso \textbf{superior al 20\%}.

%     \item \textit{Tiempo medio de retraso en la salida}:
    
%     Por otra parte y relativo al retraso en la salida, este indicador se activará siempre que 
% \end{enumerate}


% const ALERT_THRESHOLDS = {
%   regularity_min:  80,   // Regularity_value < 80
%   delay_pct_max:   20,   // Percentage_sims_delay_departure > 20%
%   delay_secs_max: 300,   // Mean_Minutes_departure_delay > 300s (5 min)
%   occupation_max:  60,   // Mean_Occupation_Percentage > 60%
% };



\section{Validación operativa y comparativa de soluciones}\label{sec:validation}



\section{Experimentación paramétrica}\label{sec:experiments}



